\chapterstart{Связанные работы и выбор подходов}
\label{sec:review}
Рассматриваемая задача объединяет три вопроса: как построить управляющий закон, как соблюсти ограничения и как корректно сравнить результаты. Литература по каждому из них даёт свои основания для экспериментального протокола. Ниже обсуждается не общий рейтинг алгоритмов, а то, какие свойства подходов существенны именно для подъёма маятника на короткой направляющей.

\subsection{Классическое управление как содержательная база сравнения}
Линейно-квадратичный регулятор (Linear Quadratic Regulator, LQR) выбирает линейную обратную связь по отклонению от равновесия, а его коэффициенты определяются из квадратичного штрафа за отклонение состояния и за усилие управления. Для нелинейного маятника такой регулятор применяют к линеаризации около верхнего положения: он решает задачу локальной стабилизации и сам по себе подъёма снизу не задаёт. Дополняет его траекторная оптимизация — сначала строится номинальное движение, затем регулятор компенсирует отклонения от него. Оба метода изложены Р.~Тедрейком~\cite{tedrake_lqr,tedrake_traj}.

Разделение подъёма и стабилизации вытекает из нелинейности задачи. Линеаризация описывает малые отклонения от выбранного движения, и её точность не распространяется автоматически на нижнее положение маятника, поэтому траекторный регулятор использует изменяющееся во времени опорное движение, а стационарный LQR предназначен для локальной стабилизации около верхнего равновесия. В реализованном здесь контроллере переключение на него выполняется по окончании запланированной траектории, то есть по времени, а не по проверке близости состояния к равновесию. Сопоставлять такой подход с обучаемой политикой нужно по полной задаче подъёма и удержания, сохраняя при этом сведения о том, как построена опорная траектория.

В проекте классическая процедура включает заранее сохранённую траекторию подъёма, слежение с обратной связью и переход к LQR. Она нужна не только как техническая проверка симулятора: если известная модель уже допускает эффективное классическое решение, успех обучаемого метода следует сопоставлять именно с ним, а не только с бездействием. Одна фиксированная траектория, впрочем, не обязана подходить всем расширенным стартам, и это учитывается при интерпретации результатов Главы~\ref{sec:robust}.

Исходный лабораторный проект \texttt{robotics-laboratory/cart-pole} предоставляет модель и основу классического управления~\cite{laboratory}. Его вход — ускорение тележки, а не сила двигателя, и уже поэтому совпадение названия CartPole с другими средами не делает результаты сопоставимыми: проверять нужно уравнения, начальную область, длину направляющей и критерий успеха. Настоящая работа сохраняет лабораторную динамику и строит вокруг неё общий экспериментальный протокол.

\subsection{Четыре подхода к обучению управляющего закона}
Все рассматриваемые методы опираются на нейросетевую оценку ценности действия — ожидаемой будущей суммы наград при его выборе и последующем поведении. Такая оценка позволяет учитывать отложенный результат: кратковременное движение тележки от центра может оказаться полезным, если впоследствии оно позволяет поднять маятник. Ценность не равна немедленной награде и сама является приближением, которое меняется по ходу обучения.

\textbf{DQN} (Deep Q-Network) обучает оценки для дискретного набора действий. Повторное использование переходов из буфера опыта и медленно меняющаяся целевая сеть снижают зависимость обучения от последовательности соседних наблюдений~\cite{mnih}. В нашей задаче DQN выбирает один из девяти запросов ускорения — это простой дискретный подход, но набор доступных ему запросов заведомо уже, чем у непрерывных методов. Общая физическая среда, как видно уже здесь, не означает одинакового класса возможных управляющих законов.

\textbf{DDPG} (Deep Deterministic Policy Gradient) предназначен для непрерывных действий. Сеть управления, \emph{актор} (actor), выдаёт действие, а сеть оценки, \emph{критик} (critic), предсказывает его ценность; актор меняется в направлении роста предсказанной ценности, а исследование обеспечивается добавленным шумом~\cite{lillicrap}. DDPG выбран как базовый непрерывный метод. Его зависимость от качества критика делает особенно полезной проверку разных seed, хотя сама по себе архитектура ещё не объясняет неуспех конкретного обучения.

\textbf{SAC} (Soft Actor-Critic) сочетает рост ожидаемого возврата с энтропией — мерой разнообразия распределения действий. Стохастическая политика поддерживает исследование, а две оценки критиков используются для более осторожной цели обновления~\cite{haarnoja_sac}; вариант с автоматической настройкой веса энтропии описан в развитии метода~\cite{haarnoja_apps}. Для маятника SAC интересен как способ находить последовательность подъёма без заранее заданной траектории. Энтропия при этом остаётся свойством распределения действий и ограничением на положение тележки не является, так что заменить защиту она не может.

\textbf{TQC} (Truncated Quantile Critics) сохраняет энтропийного актора, но критики предсказывают квантили распределения возврата вместо одного среднего, и при построении обучающей цели часть наиболее высоких прогнозов отбрасывается~\cite{kuznetsov}. Это механизм контроля переоценивания, а не встроенная проверка физических границ. Сопоставление SAC и TQC показывает результат двух полных процедур; без отдельного эксперимента с выключенным усечением оно не устанавливает, что возможное преимущество вызвано именно квантилями.

Четыре метода вместе покрывают несколько разных способов построения политики: дискретный выбор запросов, детерминированное непрерывное управление, стохастическое управление с энтропийным критерием и его развитие с распределительной оценкой возврата. Исчерпывающим обзором современных алгоритмов такой набор, конечно, не является — он выбран для содержательного сравнения при доступном бюджете. Добавление новых методов без сопоставимого подбора параметров само по себе не сделало бы выводы сильнее.

Различия алгоритмов не следует путать с доказанными причинами экспериментальных исходов, и дальше это правило применяется последовательно. Успешность TQC не доказывает пользу усечения квантилей, пока нет сравнения с отключённым усечением. Неудача DQN не устанавливает, что девяти действий недостаточно: одновременно влияют обучение оценки ценности, исследование среды и бюджет переходов. Поэтому результаты интерпретируются на уровне полных процедур; об отдельных механизмах эти данные не говорят.

Использованы реализации Stable-Baselines3 и SB3-Contrib~\cite{raffin,contrib}: библиотечные алгоритмы не переписывались и авторскими здесь не выдаются, а разработанная часть проекта описана в Главе~\ref{sec:design}. Общие небольшие сети и одинаковый бюджет взаимодействий делают сравнение обозримым, но индивидуально оптимальных настроек каждому методу не гарантируют. Выводить из полученного порядка результатов общий рейтинг алгоритмов на всех задачах управления было бы ошибкой.

\subsection{Защитный фильтр и достоверность сравнения}
Предиктивные защитные фильтры отделяют предложение контроллера от проверки его допустимости. У Wabersich и Zeilinger фильтр использует прогноз ограниченной нелинейной системы и при необходимости корректирует вход~\cite{wabersich}. В нашем случае реализована существенно более узкая процедура — аналитический прогноз координаты и скорости тележки при постоянном ускорении. Воспроизведением общего нелинейного алгоритма из этой статьи она не является и гарантий для неопределённой динамики автоматически не наследует.

Safe-control-gym объединяет оценку качества и безопасности классических и обучаемых регуляторов, в том числе при возмущениях~\cite{yuan}. Для настоящей работы важен оттуда принцип раздельных метрик: агент может достичь цели, нарушив ограничение, либо сохранить ограничение, не достигнув цели. Шумы, задержки и ошибки параметров в нашем протоколе не испытывались, поэтому численные результаты здесь с этим набором задач напрямую не сопоставляются — он служит методологическим ориентиром.

Henderson и соавторы показывают, насколько выводы о глубоком обучении с подкреплением зависят от случайности и деталей процедуры~\cite{henderson}. Отсюда прямо следуют три решения протокола: сохранять каждый запуск, оценивать последние веса отдельно от лучших и проверять выбранную настройку на новых seed. Три значения seed не дают точной оценки распределения качества будущих обучений, но позволяют увидеть крупную неоднородность, которую единственный успешный пример скрыл бы полностью.

\subsection{Позиционирование настоящего исследования}
Из рассмотренных работ следуют три требования к сравнению: определить модель и критерии явно, поскольку одинаковое название задачи не гарантирует одинаковой динамики; измерять соблюдение ограничений отдельным механизмом и отдельными показателями, поскольку высокая награда сертификатом безопасности не является; и не считать результат единственного обучения оценкой повторяемости процедуры.

Настоящий проект реализует эти требования в лабораторной модели с ускорительным входом. Его результат — не универсальное преимущество обучения с подкреплением над классическим управлением, а воспроизводимое сопоставление конкретных процедур вместе с диагностикой их ограничений. Общая модель делает физические результаты сравнимыми, а разделение отбора, новых запусков и новых стартов позволяет каждый раз сказать, на какие именно данные опирается вывод. В следующей главе эти принципы переводятся в уравнения и измеряемые критерии.
